\documentclass{article}
\usepackage{amsmath,amssymb,amsthm}
\usepackage{algorithm}
\usepackage{algpseudocode}
\usepackage{enumitem}
\usepackage{hyperref}
\newtheorem{theorem}{Theorem}
\newtheorem{lemma}{Lemma}
\newtheorem{assumption}{Assumption}
\newcommand{\EE}{\mathbb{E}}
\newcommand{\R}{\mathbb{R}}
\newcommand{\norm}[1]{\left\|#1\right\|}
\newcommand{\e}{\mathbf{e}}
\begin{document}

%=====================================================================
\section*{Randomized Coordinate Descent with Importance Sampling}
%=====================================================================

\section*{Mathematical Formulation}
Let $f:\R^{d}\rightarrow\R$ be differentiable and possibly non‑convex.
For each coordinate $i\in\{1,\dots ,d\}$ we are given a probability
$p_i>0$ with $\sum_{i=1}^{d}p_i=1$.
Denote by $\e_i\in\R^{d}$ the $i$‑th canonical unit vector.
At iteration $t$ we draw a random index $i_t\in\{1,\dots ,d\}$ according to
the distribution $\{p_i\}_{i=1}^{d}$ and perform the update  

\[
\boxed{
w_{t+1}=w_{t}-\frac{\eta}{p_{i_t}}\,
\bigl[\nabla f(w_t)\bigr]_{i_t}\,\e_{i_t}}
\tag{1}
\]

where $\eta>0$ is a stepsize (learning rate) common to all coordinates.
The factor $1/p_{i_t}$ guarantees that the stochastic update is an
unbiased estimator of a scaled full gradient (see the proof).

\section*{Pseudocode}
\begin{algorithm}[H]
\caption{Randomized Coordinate Descent (RCD) with importance sampling}
\begin{algorithmic}[1]
\Require Initial point $w_0\in\R^{d}$, stepsize $\eta>0$, probabilities
$\{p_i\}_{i=1}^{d}$ with $p_i>0$, $\sum_i p_i=1$,
maximum iterations $T$.
\For{$t=0,\dots ,T-1$}
    \State Sample $i_t\sim\{p_1,\dots ,p_d\}$.
    \State Compute the $i_t$‑th partial derivative 
          $g_{t}\gets [\nabla f(w_t)]_{i_t}$.
    \State Update the selected coordinate
          \[
          w_{t+1}\gets w_t-\frac{\eta}{p_{i_t}}\,g_{t}\,\e_{i_t}.
          \]
\EndFor
\State \Return $w_T$ (or any intermediate iterate).
\end{algorithmic}
\end{algorithm}

\section*{Dry Run Example}
Consider the one‑dimensional function $f(w)=(w-3)^2$.
Here $d=1$, $p_1=1$, $\e_1=1$, and $\nabla f(w)=2(w-3)$.
Take $w_0=0$ and stepsize $\eta=0.1$.
The update (1) reduces to the ordinary gradient step
$w_{t+1}=w_t-\eta\,\nabla f(w_t)$.

\begin{center}
\begin{tabular}{c|c|c|c}
$t$ & $w_t$ & $g_t=\nabla f(w_t)$ & $f(w_t)$ \\ \hline
0 & $0$ & $2(0-3)=-6$ & $9$ \\
1 & $0-0.1(-6)=0.6$ & $2(0.6-3)=-4.8$ & $(0.6-3)^2=5.76$ \\
2 & $0.6-0.1(-4.8)=1.08$ & $2(1.08-3)=-3.84$ & $(1.08-3)^2=3.6864$ \\
3 & $1.08-0.1(-3.84)=1.464$ & $2(1.464-3)=-3.072$ & $(1.464-3)^2=2.3660$
\end{tabular}
\end{center}
The objective value decreases at each iteration, illustrating the
mechanics of the algorithm.

%=====================================================================
\section*{Assumptions}
%=====================================================================
\begin{assumption}[Coordinate‑wise $L_i$‑smoothness]\label{ass:smooth}
For each coordinate $i\in\{1,\dots ,d\}$ there exists $L_i>0$ such that
for all $x\in\R^{d}$ and any scalar $\alpha\in\R$,
\[
\bigl|\,[\nabla f(x+\alpha \e_i)]_i-[\nabla f(x)]_i\,\bigr|
\le L_i|\alpha|.
\tag{2}
\]
Equivalently,
\[
f(x-\alpha \e_i)\le f(x)-\alpha [\nabla f(x)]_i
+\frac{L_i}{2}\alpha^2 .
\tag{3}
\]
\end{assumption}

\begin{assumption}[Lower boundedness]\label{ass:lb}
$f$ is bounded below, i.e.\ there exists $f_{\inf}>-\infty$ such that
$f(w)\ge f_{\inf}$ for all $w\in\R^{d}$.
\end{assumption}

\begin{assumption}[Stepsize]\label{ass:stepsize}
The stepsize satisfies
\[
0<\eta\le \min_{i\in[d]}\frac{p_i}{L_i}.
\tag{4}
\]
\end{assumption}

%=====================================================================
\section*{Main Convergence Theorem}
%=====================================================================
\begin{theorem}\label{thm:main}
Assume \ref{ass:smooth}--\ref{ass:stepsize} hold.
Define the (random) iterate $w_t$ generated by (1) and let
$w^{*}\in\arg\min f$ (any global minimiser, possibly non‑unique).
Then for any integer $T\ge 1$,
\[
\frac{1}{T}\sum_{t=0}^{T-1}\EE\!\left[\,
\bigl\|\nabla f(w_t)\bigr\|_{2}^{2}\,\right]
\;\le\;
\frac{2\bigl(f(w_0)-f_{\inf}\bigr)}{\eta\,T}.
\tag{5}
\]
Consequently, there exists an iteration $t\in\{0,\dots ,T-1\}$ such that
\[
\EE\!\bigl[\|\nabla f(w_t)\|_{2}^{2}\bigr]
\le \frac{2\bigl(f(w_0)-f_{\inf}\bigr)}{\eta\,T}.
\]
Thus the algorithm attains the usual $O(1/T)$ rate for the
expected squared gradient norm, even in the non‑convex setting.
\end{theorem}

%=====================================================================
\section*{Theoretical Properties}
%=====================================================================
\begin{enumerate}[label=\arabic*)]
\item \textbf{Unbiasedness of the stochastic direction.}
    From (1) and the sampling distribution we have
    \[
    \EE_{i_t}\!\bigl[\,\tfrac{1}{p_{i_t}}[\nabla f(w_t)]_{i_t}\e_{i_t}\mid w_t\bigr]
    =\sum_{i=1}^{d}p_i\,\frac{1}{p_i}[\nabla f(w_t)]_i\e_i
    =\nabla f(w_t).
    \tag{6}
    \]
    Hence the expected update equals a full gradient step with stepsize
    $\eta$.

\item \textbf{Stability w.r.t. the sampling distribution.}
    The bound (5) depends on the probabilities only through the
    admissible stepsize (4). Choosing larger $p_i$ for coordinates with
    large $L_i$ permits a larger $\eta$, yielding faster progress.

\item \textbf{Comparison with uniform coordinate descent.}
    If $p_i=1/d$ for all $i$, (4) reduces to $\eta\le \frac{1}{d\max_i L_i}$,
    reproducing the classic $O(d)$ slowdown of uniform sampling.
    Importance sampling can therefore improve the constant factor
    dramatically when $L_i$ vary across coordinates.

\item \textbf{Extension to mini‑batch coordinate updates.}
    Updating a set $S_t\subseteq[d]$ of coordinates with probabilities
    $p_i$ and scaling each update by $1/p_i$ leads to the same
    expectation (6) and the proof carries over with $L_i$ replaced by
    the smoothness constant of the selected block.
\end{enumerate}

%=====================================================================
\section*{Preliminaries (Definitions and Lemmas)}
%=====================================================================
\begin{lemma}[Descent Lemma for a single coordinate]\label{lem:descent}
Under Assumption \ref{ass:smooth}, for any $x\in\R^{d}$, any $i\in[d]$ and any
$\alpha\in\R$,
\[
f(x-\alpha \e_i)\le f(x)-\alpha [\nabla f(x)]_i
+\frac{L_i}{2}\alpha^{2}.
\tag{7}
\]
\end{lemma}
\begin{proof}
Inequality (7) is exactly the second part of Assumption
\ref{ass:smooth}. ∎
\end{proof}

%=====================================================================
\section*{Main Proof (Step‑by‑Step Derivation)}
%=====================================================================
We prove Theorem \ref{thm:main}. Throughout the proof,
expectations are taken w.r.t. the random choice of $i_t$ conditioned
on the past iterates.

\noindent\textbf{Step 1 (Apply the descent lemma).}
For the random index $i_t$ we set $\alpha=\dfrac{\eta}{p_{i_t}}
[\nabla f(w_t)]_{i_t}$ in Lemma \ref{lem:descent} and obtain
\[
f(w_{t+1})
\le f(w_t)-\frac{\eta}{p_{i_t}}[\nabla f(w_t)]_{i_t}^{2}
+\frac{L_{i_t}}{2}\Bigl(\frac{\eta}{p_{i_t}}\Bigr)^{2}
[\nabla f(w_t)]_{i_t}^{2}.
\tag{8}
\]
\textbf{[Step 1]}

\noindent\textbf{Step 2 (Take conditional expectation).}
Conditioning on $w_t$ and using $\EE_{i_t}[\,\cdot\,|w_t]$ gives
\[
\EE_{i_t}[f(w_{t+1})\mid w_t]
\le f(w_t)-\eta\sum_{i=1}^{d}[\nabla f(w_t)]_{i}^{2}
+\frac{\eta^{2}}{2}\sum_{i=1}^{d}\frac{L_i}{p_i}
[\nabla f(w_t)]_{i}^{2}.
\tag{9}
\]
\textbf{[Step 2]}

\noindent\textbf{Step 3 (Factor the gradient norm).}
Define the weighted constant
\[
\gamma_i\;:=\;\eta-\frac{\eta^{2}L_i}{2p_i}.
\tag{10}
\]
With this notation (9) rewrites as
\[
\EE_{i_t}[f(w_{t+1})\mid w_t]
\le f(w_t)-\sum_{i=1}^{d}\gamma_i\,
[\nabla f(w_t)]_{i}^{2}.
\tag{11}
\]
\textbf{[Step 3]}

\noindent\textbf{Step 4 (Use the stepsize condition).}
Assumption \ref{ass:stepsize} implies $\eta\le p_i/L_i$,
hence $\eta L_i/(2p_i)\le \tfrac12$ and consequently
\[
\gamma_i\;=\;\eta\Bigl(1-\frac{\eta L_i}{2p_i}\Bigr)
\;\ge\;\frac{\eta}{2},\qquad\forall i\in[d].
\tag{12}
\]
\textbf{[Step 4]}

\noindent\textbf{Step 5 (Obtain a uniform descent bound).}
Plugging (12) into (11) yields
\[
\EE_{i_t}[f(w_{t+1})\mid w_t]
\le f(w_t)-\frac{\eta}{2}\,\bigl\|\nabla f(w_t)\bigr\|_{2}^{2}.
\tag{13}
\]
\textbf{[Step 5]}

\noindent\textbf{Step 6 (Take full expectation).}
Denote $\EE_t[\cdot]=\EE[\cdot\mid w_0,\dots ,w_t]$.
Taking expectation of (13) with respect to all randomness up to
iteration $t$ gives the recursion
\[
\EE\!\bigl[f(w_{t+1})\bigr]
\le \EE\!\bigl[f(w_{t})\bigr]
-\frac{\eta}{2}\,
\EE\!\bigl[\|\nabla f(w_t)\|_{2}^{2}\bigr].
\tag{14}
\]
\textbf{[Step 6]}

\noindent\textbf{Step 7 (Sum the recursion).}
Summing (14) for $t=0,\dots ,T-1$ and telescoping,
\[
\EE\!\bigl[f(w_T)\bigr]
\le f(w_0)-\frac{\eta}{2}\sum_{t=0}^{T-1}
\EE\!\bigl[\|\nabla f(w_t)\|_{2}^{2}\bigr].
\tag{15}
\]
\textbf{[Step 7]}

\noindent\textbf{Step 8 (Use lower boundedness).}
By Assumption \ref{ass:lb}, $\EE[f(w_T)]\ge f_{\inf}$, thus
\[
\frac{\eta}{2}\sum_{t=0}^{T-1}
\EE\!\bigl[\|\nabla f(w_t)\|_{2}^{2}\bigr]
\le f(w_0)-f_{\inf}.
\tag{16}
\]
\textbf{[Step 8]}

\noindent\textbf{Step 9 (Derive the average bound).}
Dividing (16) by $T$ and rearranging gives exactly the claim (5):
\[
\frac{1}{T}\sum_{t=0}^{T-1}
\EE\!\bigl[\|\nabla f(w_t)\|_{2}^{2}\bigr]
\le \frac{2\bigl(f(w_0)-f_{\inf}\bigr)}{\eta\,T}.
\tag{17}
\]
\textbf{[Step 9]}

\noindent\textbf{Step 10 (Existence of a good iterate).}
Since the average of non‑negative numbers is bounded by the RHS,
there must exist at least one index $t\in\{0,\dots ,T-1\}$ such that
\[
\EE\!\bigl[\|\nabla f(w_t)\|_{2}^{2}\bigr]
\le \frac{2\bigl(f(w_0)-f_{\inf}\bigr)}{\eta\,T},
\]
which completes the proof. ∎

%=====================================================================
\section*{Rate Derivation (With Constants)}
%=====================================================================
The bound (5) can be expressed in terms of the
\emph{average smoothness constant}
\[
\bar L\;:=\;\sum_{i=1}^{d}p_i L_i,
\qquad
\text{and the minimal sampling probability }p_{\min}:=\min_i p_i .
\]
From (4) we may safely choose
\[
\eta = \frac{p_{\min}}{\bar L}\; \le\; \min_i\frac{p_i}{L_i}.
\]
Plugging this choice into (5) yields
\[
\frac{1}{T}\sum_{t=0}^{T-1}
\EE\!\bigl[\|\nabla f(w_t)\|_{2}^{2}\bigr]
\;\le\;
\frac{2\bar L}{p_{\min}}\,
\frac{f(w_0)-f_{\inf}}{T}.
\tag{18}
\]
Hence the convergence constant scales linearly with the ratio
$\bar L/p_{\min}$; importance sampling reduces this ratio compared
to uniform sampling ($p_i=1/d$) whenever the smoothness constants $L_i$
are heterogeneous.

%=====================================================================
\section*{Special Cases}
%=====================================================================
\begin{enumerate}[label=\alph*)]
\item \textbf{Uniform sampling.}
    Setting $p_i=1/d$ for all $i$ gives $\eta\le \frac{1}{d\max_i L_i}$
    and the bound (18) reduces to the classic $O(d\max_i L_i/T)$ rate.

\item \textbf{Exact coordinate‑wise Lipschitz constants.}
    If $p_i\propto L_i^{-1}$ (i.e. $p_i = \frac{L_i^{-1}}{\sum_j L_j^{-1}}$),
    then $\bar L = \bigl(\sum_j L_j^{-1}\bigr)^{-1}$ and $p_{\min}
    = \frac{1}{L_{\max}\sum_j L_j^{-1}}$.
    Consequently the factor $\bar L/p_{\min}=L_{\max}$, which is the
    best possible constant achievable by any coordinate‑wise scheme.

\item \textbf{Block coordinate descent.}
    The proof extends verbatim to blocks $B\subseteq[d]$ with a block
    smoothness constant $L_B$ and sampling probability $p_B$,
    provided the stepsize satisfies $\eta\le p_B/L_B$ for every block.
\end{enumerate}

%=====================================================================
\section*{Discussion}
%=====================================================================
The presented analysis shows that a simple importance‑sampling
strategy—weighting the update by $1/p_i$—preserves the unbiasedness of
the stochastic direction and leads to a clean $O(1/T)$ decrease of the
expected squared gradient norm, despite the non‑convex nature of $f$.
The only restriction on the stepsize is the per‑coordinate bound
(4), which is milder than the global bound required by uniform
coordinate descent. Consequently, the algorithm can be significantly
faster when the smoothness constants vary widely across coordinates.
The proof technique mirrors that of stochastic gradient descent
(see the reference examples) but adapts the smoothness inequality to
the coordinate‑wise setting and carefully tracks the $1/p_i$ factor
through the expectation calculations.

\end{document}